%!TEX root = thesis.tex

%:-------------------------- Preamble -----------------------

% Three languages are supported, which will be reflected in the logo on the front page. Pass the appropriate language
% specified as a class option to uit-thesis. Passing any other languages supported by babel will result in the default
% language on the frontpage. If no language is passed, the default is selected.
%  - USenglish (default)
%  - norsk
%  - samin
% The frontpage comes in two variants, Master's thesis and PhD. Master is default, use classoption 'phd' for the PhD version.
\documentclass[USenglish]{uit-thesis}

% Lorem ipsum
\usepackage{lipsum}

\makeglossaries

% Add external glossaryentries
\loadglsentries{acronyms}
\newacronym{api}{API}{application programming interface}\glsunset{api}
\newacronym{2api}{2API}{application programming interface}
\newacronym{d3}{D3}{Data-Driven Documents}
\newacronym{html5}{HTML5}{version 5 of the HyperText Markup Language standard}
\newglossaryentry{thesis}
{
  name=thesis,
  description={is a document submitted in support of candidature for an
    academic degree or professional qualification presenting the author's
    research and findings
    },
}
\newglossaryentry{lage}
{
  name={long ass glossary entry},
  description={is a long ass entry with a lot of text describing the properties of the glossary entry. Hopefully this spans some lines now.
  },
}


\newcommand{\listdefinitionname}{My list of definitions}
\newlistof{definition}{def}{\listdefinitionname}
\newcommand{\definition}[1]{%
  \refstepcounter{definition}%
  \par\noindent\textbf{The Definition~\thedefinition. #1}%
  \addcontentsline{def}{definition}
    {\protect\numberline{\thechapter.\thedefinition}#1}\par%
}

\begin{document}

%:-------------------------- Frontpage ------------------------

\title{Title of the master thesis}
\subtitle{Subtitle}			% Optional
\author{Name of author}
\thesisfaculty{Faculty of Science and Technology \\ Department of Computer Science}
\thesisprogramme{Master thesis in [study programme] … Month 20xx}
%\ThesisFrontpageImage{example_image.jpg}	% Optional

\maketitle

%:-------------------------- Frontmatter -----------------------
\frontmatter

\begin{dedication}
To Leslie.

Fuck you very much.
\end{dedication}

\begin{epigraph}
\epigraphitem{Simplicity is prerequisite for reliability.}{Edsger Dijkstra}
\epigraphitem{Beware of bugs in the above code;\\I have only proved it correct, not tried it.}{Donald Knuth}
\end{epigraph}

\begin{abstract}
\lipsum[2-3]
\end{abstract}

\begin{acknowledge}
\lipsum[4-8]
\end{acknowledge}

\tableofcontents

\listofdefinition

%:-------------------------- Mainmatter -----------------------
\mainmatter

\chapter{Introduction}
\lipsum[1]

\section{First section}
\lipsum[2-6]

\section{Second section}
\lipsum[6]

\begin{itemize}
\item The first item is a glorious item! It may span several lines and even several paragraphs. Hopefully this spans some lines
now. Sed feugiat. Cum sociis natoque penatibus et magnis dis parturient montes, nascetur ridiculus mus. Et requiem morbi et nunc. Aliquam dolor odio.

This is the second paragraph of that first item. Gorgeous, gorgeous, gorgeous! Nulla non mauris vitae wisi pisuere convallis. Sed eu nulla nec eros scelerisque pharetra. Nullam varius.
\item The second item is a peculiar and rather scrumptious item, I would say. Quite, quite interesting.
\end{itemize}

\lipsum[7]

\begin{enumerate}
\item The first item is a glorious item! It may span several lines and even several paragraphs. Hopefully this spans some lines
now. Sed feugiat. Cum sociis natoque penatibus et magnis dis parturient montes, nascetur ridiculus mus. Et requiem morbi et nunc. Aliquam dolor odio.

This is the second paragraph of that first item. Gorgeous, gorgeous, gorgeous! Nulla non mauris vitae wisi pisuere convallis. Sed eu nulla nec eros scelerisque pharetra. Nullam varius.
\item The second item is a peculiar and rather scrumptious item, I would say. Quite, quite interesting.
\end{enumerate}

\newpage

\section{Third section}
\lipsum[6-8]

\begin{description}
\item[Entry A] with definition A.
\item[Entry B] with definitioin B.
\item[Entry C] with definition C.
\end{description}

\section{Fourth section}
\lipsum[6-7]

\subsection{A subsection}
\lipsum[8-10]

\subsubsection{A subsubsection}
\lipsum[10-11]

\paragraph{A paragraph}
\lipsum[12]

\section{Fifth section}
\lipsum[6-7]

\subsection{A subsection}
\lipsum[8-9]


We can use the \ac{api} to \ac{2api} do stuff, and write about what we did in a \gls{thesis}!

This is some stuff, {\sc smallcaps {\em smallcapsemphasized}} {\em regularemphasized}

\Gls{lage}: a test glossary entry.

If the acronym \ac{uit} is displayed, then loadglsentries works.

It is fun to use modern \upsc{OpenMP} technology!\footnote{This is a snarky
footnote. Words and etc. Semantic web technologies are technologies that enable
semantification of the Web as we know it today. Hopefully this spans some lines
now.}

It is fun to use \emph{modern \upsc{OpenMP}} technology! And it is fun to use \ac{d3} and \ac{html5}.

Referencing figure \ref{fig:ex} to test link.\footnote{This is another
footnote.}

\definition{Some other definition}

\begin{figure}
\centering
\includegraphics[scale=0.1]{example_image.jpg}
\caption{Figure link should point to top of figure.}
\label{fig:ex}
\end{figure}

\chapter{Design}
\lipsum[8-10]

\definition{Another definition}

\begin{table}
\centering
\begin{tabular}{|l|l|}
\hline
Content left & Content right\\
\hline
\end{tabular}
\caption{A table}
\end{table}

\begin{table}
\centering
\begin{tabular}{|l|l|}
\hline
Content left & Content right\\
\hline
\end{tabular}
\caption{Another table}
\end{table}

\newpage

\begin{lstlisting}[frame=single,caption={Small C program},language=C]
#include "stdio.h"
#define e 3
#define g (e/e)
#define h ((g+e)/2)
#define f (e-g-h)
#define j (e*e-g)
#define k (j-h)
#define l(x) tab2[x]/h
#define m(n,a) ((n&(a))==(a))

long tab1[]={ 989L,5L,26L,0L,88319L,123L,0L,9367L };
int tab2[]={ 4,6,10,14,22,26,34,38,46,58,62,74,82,86 };

main(m1,s) char *s; {
  int a,b,c,d,o[k],n=(int)s;
  if(m1==1){ char b[2*j+f-g]; main(l(h+e)+h+e,b);
    printf(b); }
  else switch(m1-=h){
    case f:
      a=(b=(c=(d=g)<<g)<<g)<<g;
      return(m(n,a|c)|m(n,b)|m(n,a|d)|m(n,c|d));
    case h:
      for(a=f;a<j;++a)
        if(tab1[a]&&!(tab1[a]%((long)l(n))))
          return(a);
    case g:
      if(n<h)return(g);
      if(n<j){n-=g;c='D';o[f]=h;o[g]=f;}
      else{c='\r'-'\b';n-=j-g;o[f]=o[g]=g;}
      if((b=n)>=e)
        for(b=g<<g;b<n;++b)o[b]=o[b-h]+o[b-g]+c;
      return(o[b-g]%n+k-h);
    default:
      if(m1-=e) main(m1-g+e+h,s+g); else *(s+g)=f;
      for(*s=a=f;a<e;) *s=(*s<<e)|main(h+a++,
      (char *)m1);
    }
}
\end{lstlisting}

\chapter{The single-line chapter title}
\lipsum[1-2]

\chapter{This is a long chapter title that spans two lines}
\lipsum[3-4]

\chapter{This is an even longer chapter title that spans three lines}
\lipsum[5-6]

\backmatter

\end{document}

